A partial linear space is a set system on a fixed set $V$. 
We write $\cL=(V,\cB),$ 
where $\cB$ is a set of distinct subsets of $V$, called lines. 
The members of $V \cup \cB$ are called elements. 
For two elements $x,y$, we say that $x$ is incident with $y$, written $xIy$, 
if either $x \in y$ or $y \in x.$
We require that any line has at least two points 
and that any two points are contained in at most one line. 
A {\em decomposition} of a linear space 
is a partition $\Pi= (C_1,\ldots,C_n)$ of $V \cup \cB$ 
such that each $C_i$ either is a subset of $V$ or a subset of $\cB.$
A decomposition is called {\em tactical} if for all $i,$ the incidence number 
$$
\iota(C_i,C_j) = \# \{y \in C_j, xIy\}
$$
does not depend on the choice of $x \in C_i.$
Any linear space has a tactical decomposition, 
as the discrete partition (every element is in its own class) 
is tactical. 
Let $\Aut(\cL)$ be the automorphism group of the linear space, 
which is the subgroup of 
$\Sym(V)$ which preserves incidence. 
For $\alpha \in \Aut(\cL),$ we say that 
the decomposition $\Pi$ preserves $\alpha$ if $\alpha$ fixes every class of $\Pi.$ 
For $A \le  \Aut(\cL),$ we say that $\Pi$ preserves $A$ 
if $\Pi$ preserves every element $\alpha \in A.$
Mostly, we are interested in those decompositions $\Pi$ which preserve $\Aut(\cL).$
In light of this, 
the discrete decomposition is not that interesting. 

\bigskip


Any linear space has a coarsest tactical decomposition 
that preserves its automorphism group: 
The orbit partition of the automorphism group acting on $V \cup \cB$ will do. 
Up to ordering of the classes, the coarsest tactical refinement is unique.
Computing the orbit decomposition is challenging 
as it involves computing the automorphism group. 
Computationally, there are easier ways to get to admissible decompositions. 
One way is by means of successive refinements. 
If a class $C_i$ does not have the property that $\iota(C_i,C_j)$ is well-defined 
for all $x \in C_i,$ then a refinement of $C_i$ will do. 
The coarsest refinement of $C_i$ has the property that 
if $C_i$ preserves some group $A$ then the refinement will do, too.
This shows that there is an algorithm to compute 
a tactical decompostion of any given linear space $\cP.$ 
Simply start with the decomposition of two classes, 
one the set of points and one the set of blocks, and refine. 
The output may or may not be equal to the decomposition 
arising from the orbit partition of $\Aut(\cL).$

\bigskip

Let us consider the opposite question. 
Given a tactical decomposition, 
can we classify the linear spaces which admit a given tactical decomposition.
The \verb'-geometry_builder' option can help. 
Table~\ref{tab:geometry:builder} shows the options for the geometry builder.
\orbitertableofcommands{geometry_builder}


\bigskip

The command 

\sourcecode{geo_10_3}

classifies the configuations $10_3$. 
It uses isomorphism tests after $4,$ $5,$ $6,$ $7,$ $8,$ $9$ and $10$ points.
The positions of the tests is defined using a set called \verb'Test_lines'. 
The set of test lines is defined using a loop command.
The command shows that 
there are exactly 10 configurations of this kind. 
One of them is the Desargues configuration. 
Four different output files can be written.
Each contains all geometries, but the file format is different.
\begin{enumerate}
\item
The option \verb'-output_to_inc_file'
writes \verb'10_3.inc'. 
The file contains the incidences in increasing order. 
The position in the incidence matrix is given.
Apart from the header and footer, each row in the file describes exactly one geometry.
The incidences are given in numeric form. 
Each incidence is the numerical position of the point/block pair in the 
incidence matrix. 
The position is the numbering of the matrix entries 
in the incidence matrix in row-major ordering, 
starting with zero for the top left entry. 
The index of the incidence in row $i$ (zero-based) 
and column $j$  is $b\cdot i + j,$ 
where $b$ is the number of blocks in the geometry.
The header consists of one row, containing the number of points, 
the number of lines, and the number of incidences.
The footer consists of two rows. 
A single $-1$ indicates the beginning of the footer.
The second row in the footer summarizes the automorphism group orders of the germetries.
Here is an example:
The file \verb'10_3.inc' lists the 10 configurations $10_3$:
\orbiteroutput{GRAPHICS/LATEX/conf_10_3_inc_file}
\item
The option \verb'-output_to_sage_file'
writes \verb'10_3.sage'. 
This file is meant to be read by Sage~\cite{SteinJoyner2005}.
\item
The option \verb'-output_to_blocks_file'
writes \verb'10_3.blocks'. 
Here is the content of the file \verb'10_3.blocks':
\orbiteroutput{GRAPHICS/LATEX/conf_10_3_blocks_file}
contains the blocks. 
Each number represents the rank of a 3-subset corresponding to a block 
in the lexicographic ordering of all 3-subsets. 
\item
The option \verb'-output_to_blocks_latex_file'
writes \verb'10_3.blocks_long'. 
The file 
\verb'10_3.blocks_long' contains a list of all blocks 
written out in a format ready for use in latex.
\end{enumerate}

\bigskip

It is possible to create graphical representations of the search tree. 
The additional option \verb'-search_tree' can be given, as shown in the following command:

\sourcecode{geo_10_3_tree}

The \verb'-search_tree' option causes Orbiter 
to create a file containing the search tree. 
The name of the file is derived from the file name given with the \verb'fname_GEO' option. 
Here, the \verb'fname_GEO' option sets the output file to \verb'10_3'. The \verb'-search_tree'
option then creates the file \verb'10_3_tree.txt'.
Once the classification is completed,
the \verb'-tree_draw' command can be used to create 
a drawing of the search tree from the file just created.
The vertex color represents the outcome of the isomorphism test. 
A green node is accepted. A red node is rejected. 
The search will continue for the green nodes only. 
A green node at the bottom of the tree corresponds 
to an isomorphism type of a geometry satisfying all the requirements. 
In the example of the configurations $10_3$, there are 10 green nodes at the very bottom 
of the search tree. 
They represent the 10 isomorphism types of configurations $10_3$.
The size of the tree can be determined by counting 
the lines of the file \verb'10_3_tree.txt' and subtracting one:
\orbiteroutput{GRAPHICS/LATEX/conf_10_3_wc}
The word count command yields the following output:
\orbiteroutput{GRAPHICS/LATEX/conf_10_3_wc_output}
This means that there are 1471 lines in the file. 
Hence the search tree has 1470 nodes.
The resulting tree is shown below
\orbiteroutput{GRAPHICS/WRAPPERS/conf_10_3_search_tree}


\bigskip


Any incidence structure defines a graph on its underlying set of points. 
The vertices of the collinearity graph are the points of the incidence structure.
Two vertices are adjacent if the incidence structure contains a block 
which contains the associated points. 
The distance between two points in the geometry 
is the distance of the associated vertices in the collinearity graph. 
The girth is the length of the shortest cycle. 
Incidence structures with high girth are often interesting.
Orbiter allows classifying incidence structures with a given girth.
More specifically, Orbiter classifies 
all incidence structures whose girth is at least a given number, 
or proves that there are none.
This is interesting because the girth can be used as a filter, 
and classifying with a given girth may be 
faster than classifying all incidence geometries.
For instance, 
the following example classifies 
all cubic graphs on 10 vertices with girth at least 5:

\sourcecode{geo_petersen}

There is a unique graph with these properties. 
It is the Petersen graph. 
Its automorphism group is $\Sym(5)$ of order $120.$

\bigskip

Let us look at the girth of configurations.
For instance, while there are $245342$ 
isomorphism classes of configurations $15_3$, 
only one of them has girth 4. 
This is the Cremona Richmond configuration 
with automorphism group $\Sym(6)$ of order 720. 
It is associated to a cubic surface.
Let us compare the classification with and without girth condition.
The following command classifies all configurations $15_3$:

\sourcecode{15_3.inc}

The next command classifies the configurations $15_3$ with girth $4.$ 

\sourcecode{geo_15_3_g4}


\bigskip


The next command classifies the configurations $40_4$ with girth $4.$ 
Exactly two configuration arise, both with a group of order 51840. 
Note the extra option \verb'-special_test_not_orderly' to speed up the search.

\sourcecode{geo_40_4_g4}

The search tree is shown below
\orbiteroutput{GRAPHICS/WRAPPERS/conf_40_4_g4_search_tree}



\bigskip

The next command classifies the configurations $63_3$ with girth $6.$ 
Exactly two configuration arise, both with a group of order 12096. 
Note the extra option \verb'-special_test_orderly' to speed up the search.

\sourcecode{geo_63_3_g6}

The configuration is also known as the split Cayley hexagon.


\bigskip

The next command classifies the Latin squares of order $6.$ 

\sourcecode{geo_LSQ6.inc}

Up to isomorphism, there are exactly 12 such objects.


\bigskip




